\documentclass[tikz,border=1cm]{standalone}
\usepackage{xcolor}
\usepackage{amsmath,amssymb}

\usepackage[utf8]{inputenc}
\usepackage[spanish]{babel}
\usepackage[
    letterpaper,
    top=1.5cm,
    bottom=2cm,
    left=2cm,
    right=2cm
]{geometry}

\usepackage{graphicx}
\usepackage{fancyhdr}
\usepackage[table,xcdraw]{xcolor}
\usepackage{tikz}
\usetikzlibrary{shapes.geometric, arrows, positioning, fit, backgrounds, calc, babel}
\usepackage{ifpdf}
\usepackage{blindtext}
\usepackage[cmex10]{amsmath}
\usepackage{amssymb}
\usepackage{amsfonts}
\usepackage{float}
\usepackage{siunitx}
\usepackage{listings}
\usepackage{hyperref}
\usepackage{subcaption}
\usepackage{booktabs}
\usepackage{url}
\usepackage{wrapfig}
\usepackage{titlesec}
\usepackage[export]{adjustbox}
\usepackage{imakeidx}
\usepackage{parskip}
\usepackage[numbers]{natbib}

\usepackage{ragged2e}
\usepackage{setspace}
\usepackage{array}
\usepackage{longtable}

% Definición de un nuevo tipo de columna para evitar "Underfull \hbox" en tablas
\newcolumntype{L}[1]{>{\RaggedRight\arraybackslash}p{#1}}

\hypersetup{colorlinks=false,bookmarksopen=true,linkbordercolor={1 1 1}}

\newcommand{\rfigura}[1]{Fig.\ref{#1}}
\newcommand{\rtabla}[1]{Tabla \ref{#1}}

% Datos del informe
\newcommand{\Curso}{IEE2913 --- Diseño Eléctrico (Capstone)}
\newcommand{\TituloInforme}{Informe de Avance 1 (5\%)}
\newcommand{\NumeroGrupo}{10}
\newcommand{\NombreProyecto}{SupaClock: Wearable Biométrico Modular}
\newcommand{\Integrantes}{Tomás Avendaño, Benjamín Sepúlveda, Pablo Uribe}
\newcommand{\FechaEntrega}{7 de abril de 2026}

\lstset{
  basicstyle=\ttfamily\small,
  breaklines=true,
  breakatwhitespace=true,
  postbreak=\mbox{\textcolor{red}{$\hookrightarrow$}\space},
  captionpos=b,
  frame=single,
  numbers=left,
  numberstyle=\tiny\color{gray},
  language=C,
  keywordstyle=\color{blue}\bfseries,
  commentstyle=\color{gray}\itshape,
  stringstyle=\color{red!70!black}
}


\begin{document}
\begin{tikzpicture}[
            node distance=0.8cm and 1.2cm,
            block/.style={rectangle, draw=black, thick, fill=white, text width=2.5cm, align=center, minimum height=1.2cm, font=\footnotesize},
            arrow/.style={thick, ->, >=stealth},
            stage/.style={font=\scriptsize\bfseries, blue!80!black}
            ]
            
            % Blocks
            \node[block, fill=green!15] (sensor) {Sensores\\(AD8232 / MAX30102)};
            \node[block, fill=blue!15, right=0.8cm of sensor] (adc) {MCU (ESP32-C3)\\+ Buffer Circular};
            \node[block, fill=cyan!15, right=0.8cm of adc] (ble) {Gateway BLE\\MTU 247 Bytes};
            \node[block, fill=orange!20, right=0.8cm of ble] (app) {App Móvil\\(Flutter)};
            \node[block, fill=red!10, right=0.8cm of app] (cloud) {Cloud Firestore\\+ Functions};
            
            % Labels
            \node[stage, above=0.3cm of sensor] {1. Adquisición};
            \node[stage, above=0.3cm of adc] {2. Digitalización};
            \node[stage, above=0.3cm of ble] {3. Transmisión};
            \node[stage, above=0.3cm of app] {4. Recepción};
            \node[stage, above=0.3cm of cloud] {5. Procesamiento Nube};
            
            % Arrows
            \draw[arrow] (sensor) -- node[above, font=\scriptsize] {ADC / I2C} (adc);
            \draw[arrow] (adc) -- node[above, font=\scriptsize] {Notificaciones} (ble);
            \draw[arrow] (ble) -- node[above, font=\scriptsize] {GATT} (app);
            \draw[arrow] (app) -- node[above, font=\scriptsize] {REST / Wi-Fi} (cloud);
            
            % Decorator (Pan-Tompkins)
            \node[below=0.6cm of cloud, font=\scriptsize, align=center, draw=black!40, dashed, fill=red!5, inner sep=0.2cm] (algo) {\textbf{Algoritmo Pan-Tompkins}\\Filtro Pasa-Banda\\Detección Picos R};
            \draw[->, >=stealth, dashed, thick] (cloud) -- (algo);
            
        \end{tikzpicture}
\end{document}
